\chapter{Conclusion et remerciements}

\section{Synthèse du projet}
Ce projet \textbf{Badge Reader} a permis de mettre en œuvre les compétences acquises durant ma formation \textbf{CDA} dans un contexte professionnel concret.
J’ai conçu et développé une \textbf{application de gestion de badges} pour Colint, en utilisant une architecture moderne et des outils DevOps.
\newline

\begin{for_you}{Bilan technique et métriques}
    \textbf{Architecture implémentée} :
    \begin{itemize}
        \item \textbf{Type} : MVC Monolithique Modulaire.
        \item \textbf{Interface (Vue/Contrôleur)} : Phoenix LiveView (temps réel, sécurité serveur).
        \item \textbf{Logique \& Données (Modèle)} : Elixir et PostgreSQL (Ecto).
        \item \textbf{DevOps} : Docker, GitHub Actions, CI/CD automatisé.
        \newline
    \end{itemize}

    \textbf{Chiffres clés} :
    \begin{center}
    \begin{tabular}{|l|c|c|}
    \hline
    \textbf{Métrique}               & \textbf{Valeur} & \textbf{Objectif} \\
    \hline
    Durée de développement           & 5 mois          & 6 mois          \\
    Couverture de tests             & 92\%            & 85\%            \\
    Temps de réponse (P95)          & 180 ms          & < 500 ms        \\
    Vulnérabilités sécurité         & 0               & 0               \\
    Taux d’adoption utilisateurs    & 88\%            & 80\%            \\
    Réduction des accès non autorisés & 100\%          & 100\%           \\
    \hline
    \end{tabular}
    \end{center}

    \textbf{\newline Technologies maîtrisées} :
    \begin{itemize}
        \item \textbf{Elixir/Phoenix} : Développement full-stack, LiveView.
        \item \textbf{PostgreSQL} : Modélisation des données, requêtes optimisées.
        \item \textbf{Docker} : Containerisation, déploiement reproductible.
        \item \textbf{GitHub Actions} : CI/CD automatisé, tests et déploiement.
        \item \textbf{Sécurité} : Chiffrement, RGPD, gestion des rôles.
    \end{itemize}
\end{for_you}

\newpage

\section{Perspectives d'évolution}
Le projet \textbf{Badge Reader} peut évoluer vers une \textbf{v2.0} avec des fonctionnalités avancées, tout en conservant l’architecture actuelle pour une évolution progressive.

\begin{for_you}{Roadmap technique v2.0}
    \textbf{Fonctionnalités prévues} :
    \begin{verbatim}
    Q1 2026 : Améliorations majeures
    +-- Badges dématérialisés (iOS/Android)
    +-- Intégration avec l’intranet Colint

    Q2 2026 : Sécurité et performance
    +-- Authentification biométrique
    +-- Audit des accès (logs détaillés)
    +-- Optimisation des requêtes (P95 < 100ms)
    +-- Sauvegardes automatiques chiffrées

    Q3 2026 : Évolutions technologiques
    +-- Intégration de LiveView Native (mobile)
    +-- Tests de charge automatisés
    \end{verbatim}

    \textbf{Compétences à développer} :
    \begin{itemize}
        \item \textbf{Mobile} : LiveView Native pour les badges dématérialisés.
        \item \textbf{IA} : Détection d’anomalies (ex. : accès suspects).
        \item \textbf{Leadership} : Gestion d’équipe et revues de code.
    \end{itemize}
\end{for_you}

\section{Bilan personnel et apprentissages}
Ce projet a été une \textbf{expérience formatrice} à plusieurs niveaux :

\begin{itemize}
    \item \textbf{Technique} : Maîtrise d’Elixir/Phoenix et transition vers la programmation fonctionnelle.
    \item \textbf{Métier} : Appréhension des enjeux critiques liés à la sécurité physique et numérique.
    \newline
\end{itemize}

\begin{for_you}{Ce que j’ai appris}
    \textbf{Compétences techniques} :
    \begin{itemize}
        \item Développer une application \textbf{full-stack} avec Elixir/Phoenix.
        \item Concevoir une \textbf{base de données relationnelle} (PostgreSQL).
        \item Automatiser le déploiement avec \textbf{Docker} et \textbf{GitHub Actions}.
        \item Appliquer les \textbf{bonnes pratiques de sécurité} (RGPD, chiffrement).
        \newline
    \end{itemize}

    \textbf{Compétences transverses} :
    \begin{itemize}
        \item \textbf{Autonomie} : Résolution de problèmes techniques complexes.
        \item \textbf{Adaptabilité} : Ajuster le projet aux retours utilisateurs.
        \item \textbf{Rigueur} : Respect des deadlines et des exigences qualité.
        \item \textbf{Communication} : Présentation claire aux parties prenantes.
        \newline
    \end{itemize}

    \textbf{Difficultés surmontées} :
    \begin{itemize}
        \item \textbf{Apprentissage d’Elixir} : Langage fonctionnel nouveau pour moi.
        \item \textbf{Déploiement} : Configuration de Docker et CI/CD.
        \item \textbf{Tests} : Couverture complète des cas d’usage.
        \newline
    \end{itemize}
\end{for_you}

\section{Remerciements}
Je tiens à remercier toutes les personnes qui ont contribué à la réussite de ce projet et à mon apprentissage durant cette formation.
\newline

\begin{for_you}{Remerciements personnalisés}
    \textbf{À l’équipe de Colint : \newline}
    Pour leur \textbf{confiance} et leurs \textbf{retours constructifs}, essentiels pour affiner le produit.

    \textbf{Aux membres du jury : \newline}
    Pour leur \textbf{temps} et leur \textbf{patience}, essentiels pour passer notre diplôme.

    \textbf{À mes formateurs CDA : \newline}
    Pour m’avoir transmis les \textbf{fondamentaux techniques} et méthodologiques.

    \textbf{À la communauté Elixir : \newline}
    Pour les \textbf{ressources ouvertes} et l’entraide sur les forums.

    \textbf{À mes proches : \newline}
    Pour leur \textbf{soutien moral} durant les phases intenses du projet.
\end{for_you}

\textbf{\newline}

\begin{for_you}{Ce que je retiens}
    \textbf{Sur le plan technique} :
    \begin{itemize}
        \item L’importance d’une \textbf{architecture bien conçue} dès le départ.
        \item La puissance d’Elixir/Phoenix pour les applications \textbf{temps réel}.
        \item L’utilité des \textbf{tests automatisés} pour garantir la qualité.
        \newline
    \end{itemize}

    \textbf{Sur le plan humain} :
    \begin{itemize}
        \item La \textbf{collaboration} avec les utilisateurs est clé.
        \item La \textbf{veille technologique} est indispensable.
        \item La \textbf{persévérance} paie toujours.
        \newline
    \end{itemize}
\end{for_you}

\section{Liens utiles}

\begin{itemize}
    \item Colint.school: \url{https://colint.school/}
\end{itemize}
